\begin{conclusions}
    
    El desarrollo de este trabajo permitió predecir la eficacia del candidato vacunal Quimi-Vio 11 mediante un estudio de inmunopuente con la vacuna Quimi-Vio heptavalente, previamente aprobada en la población pediátrica. El criterio de no inferioridad se cumplió en siete de los nueve serotipos evaluados: 1, 5, 6A, 19A, 6B, 18C y 19F.
    
    Por su parte, los serotipos que no superaron el margen de no inferioridad establecido en el contraste de hipótesis (14 y 23F) coincidieron con aquellos escenarios donde la respuesta opsonofagocítica en el grupo pediátrico fue superior a la observada en la población adulta.
    
    El estudio también evidenció un elevado nivel de censura en los datos de ambos ensayos clínicos, comportamiento que fue particularmente marcado en la respuesta opsonofagocítica de la población pediátrica. 
    
    El análisis comparativo de los enfoques estadísticos demostró que, ante la presencia de una alta censura, el método de sustitución produce intervalos de confianza artificialmente más estrechos en comparación con el método de estimación por máxima verosimilitud. Para los serotipos con bajo porcentaje de censura, el comportamiento de las medias geométricas tiende a la convergencia, llegando a ser idéntico en escenarios con ausencia de censura.
    
    Estos hallazgos confirman que, ante datos censurados, la elección del método estadístico es determinante en los estudios de inmunopuente y la investigación inmunológica en general. De este modo, se evitan conclusiones sesgadas que comprometan la validez al inferir la eficacia vacunal.
\end{conclusions}
